\chapter{การทำนายและการวิเคราะห์โครงสร้าง 3 มิติของโปรตีนด้วย AlphaFold DB}

การพับตัวของโปรตีน (Protein Folding) เป็นหนึ่งในปริศนาทางชีววิทยาที่ยิ่งใหญ่ที่สุด การมาถึงของ AlphaFold ช่วยเปลี่ยนผ่านจากการทดลองในห้องปฏิบัติการที่ใช้เวลาหลายเดือน สู่การทำนายโครงสร้างระดับอะตอมได้ภายในไม่กี่นาที

\section{สถาปัตยกรรม Evoformer และการพับตัวของโครงสร้าง}

โมเดล AlphaFold ใช้สถาปัตยกรรม Evoformer ในการประมวลผลการจัดเรียงลำดับเบสหลายชุด (Multiple Sequence Alignment: MSA) ร่วมกับเมทริกซ์ระยะห่างของคู่กรดอะมิโน (Residue Pair Representations) เพื่ออนุมานความสัมพันธ์เชิงพื้นที่ใน 3 มิติ

\section{การแปลความหมายดัชนีความเชื่อมั่น pLDDT}

ค่า Predicted Local Distance Difference Test (pLDDT) มีช่วงคะแนนตั้งแต่ 0 ถึง 100 ใช้บ่งชี้ความมั่นใจของโมเดล:
\begin{itemize}
    \item \textbf{pLDDT > 90 (สีน้ำเงินเข้ม):} ความเชื่อมั่นสูงมาก โครงสร้างสายข้าง (Side Chains) มีความแม่นยำระดับอะตอม
    \item \textbf{70 < pLDDT < 90 (สีฟ้าอ่อน):} ความเชื่อมั่นระดับดี โครงสร้างแกนหลัก (Backbone) เชื่อถือได้
    \item \textbf{50 < pLDDT < 70 (สีเหลือง):} ความเชื่อมั่นต่ำ ควรระมัดระวังในการแปลผล
    \item \textbf{pLDDT < 50 (สีส้มแดง):} บริเวณที่ไม่มีโครงสร้างชัดเจน (Intrinsically Disordered Regions: IDR)
\end{itemize}

\section{การสืบค้นความคล้ายคลึงของโครงสร้างด้วย Foldseek}

เมื่อต้องการเปรียบเทียบโครงสร้างที่ทำนายได้กับคลังข้อมูล PDB ทั่วโลก Foldseek แปลงโครงสร้าง 3 มิติให้อยู่ในรูปตัวอักษร 3Di (3D Interaction Alphabet) ทำให้ค้นหาโครงสร้างที่คล้ายกันได้เร็วกว่าวิธีเดิมนับพันเท่า
